\section{Выводы}

Корпус выбран: кино и сериалы, два источника --- русская Википедия и Кино-Театр.Ру.
Документов хватает с запасом, готовые поисковые системы для сравнения есть.

Главное, что я вынес из работы: \enquote{размер документа} --- не одна величина, а три.
У статьи \enquote{Во все тяжкие} это 1152~КБ сырого HTML, 161~КБ текста и 18 тысяч слов,
причём отношение между ними непостоянно --- от 1\,\% у заглушек до 32\,\% у авторской
статьи.

Ещё два наблюдения. Статьи-заглушки формально относятся к нужной тематике, но 120 слов
текста делают их бесполезными и занижают все средние --- отсеивать их надо сразу. И
кодировки у источников разные: пока не заметишь windows-1251 на Кино-Театр.Ру, текст
оттуда выглядит нечитаемым мусором.

\subsection*{Недостатки работы}

Выборка маленькая --- 12 документов, из них 5 по второму источнику, и отбирались они не
случайно, а специально разнотипными. Средние по такой выборке неустойчивы, тем более
при большом разбросе размеров. Метод выделения текста грубый: теряется карточка фильма,
зато в текст попадают списки источников. Объём второго источника остался неизвестным.

Всё это стоит пересчитать по полному корпусу, когда он будет собран, и перейти от
средних к медианам.

\pagebreak
